\begin{table}[H]
\centering
\begin{tabular}{|p{3.5cm}|p{9.9cm}|}
\hline
\rowcolor{fptOrange}
\textbf{Category} & \textbf{Tool} \\
\hline
Project tracking & GitHub issues, milestones, and the project board in the same repository as the code \\
\hline
Version control & Git with GitHub; two long-lived branches and short-lived use-case branches \\
\hline
Communication & Discord for daily coordination, Google Meet for reviews and supervisor meetings \\
\hline
Documentation & LaTeX sources in the repository, Google Drive for shared material, Microsoft Word where a form is required \\
\hline
UI/UX design & Figma \\
\hline
Runtime and workspace & Node.js with pnpm workspaces and Turborepo task orchestration \\
\hline
Frontend & Next.js 16 App Router, React 19, Tailwind CSS 4, the shared component library, next-intl for English and Vietnamese, TanStack Query for server state, Tiptap for the forum editor \\
\hline
Backend & NestJS 11 with a global API prefix, Zod validation through nestjs-zod, a permission guard, and Socket.IO for messaging \\
\hline
Data access & Prisma 7 with the PostgreSQL driver adapter; migrations generated from the schema and applied per environment \\
\hline
Database & PostgreSQL \\
\hline
Shared contracts & \texttt{packages/shared-type}, containing the Zod schemas and the types used by both applications \\
\hline
Payments & Stripe, in test mode throughout development: Connect Express, PaymentIntent, Checkout Session, transfers, refunds, and a signature-verified webhook \\
\hline
Media storage & Cloudinary for profile and portfolio images; identity documents stored on the application's own storage and served through authenticated endpoints \\
\hline
Email & SMTP through Nodemailer for one-time verification and password reset codes \\
\hline
External AI services & A curriculum vitae analysis service and an identity verification service, both called only from the backend over HTTP with an explicit timeout \\
\hline
Testing & Jest 30 for unit and integration tests, an end-to-end test against the running API, and Postman for exploratory API checks \\
\hline
Quality tooling & ESLint 9 and Prettier, plus the TypeScript compiler in strict mode through the shared configuration package \\
\hline
Secret management & Infisical, injecting environment variables at start-up; no secret is stored in the repository \\
\hline
Containerisation & Docker with Docker Compose, providing the API and the web client; the API exposes a health check and the web client starts only after the API is healthy \\
\hline
API documentation & A maintained reference in the repository describing each endpoint, its permission, and its request and response schemas \\
\hline
\end{tabular}
\caption{Tools and Infrastructures}
\label{tab:tools}
\end{table}

The tool selection follows a single constraint: a member must be able to reproduce the whole system from a fresh clone with one command. That is why the runtime, the database, the task pipeline, and the container definition are all pinned in the repository rather than installed manually per machine. It is also why the configuration is injected through a secret manager instead of being copied between members, and why the API reference is kept in the repository rather than in a separately hosted portal.

Two capabilities are deliberately not part of the toolchain. There is no automated continuous integration service: the same quality gate is executed locally and inside the container build, and the reviewer confirms the result before approving a pull request. There is no hosted API documentation portal; the reference is a maintained document, which is sufficient for a team of five and avoids a second source of truth that can fall out of date.
